\chapter{Evaluation}
\label{chap:eval}
%%%%%%%%%%%%%%%%%%%%

This chapter evaluates the performance of the machine learning classifiers and the downstream hardware parameter optimizer. 

\section{Evaluation Setup}

\subsection{Metrics}
The segmenting classifier's quality is primarily measured via per-class classification accuracy alongside the absolute temporal detection error, expressed as the average millisecond divergence from the ground truth intention timestamp. 
The parameter optimizer is evaluated on its success in minimizing the simulation's tripartite loss: maintaining peak precision and recall whilst simultaneously minimizing latency.

\subsection{Baselines}
The optimizer evaluates hardware configuration changes linearly against the fixed threshold baselines (e.g., default firmware behaviour). The underlying detection engine evaluates success compared to raw heuristics mapping logic.

\section{Click-Intent Classifier Results}
Tasked with interpreting the analog stream, the Random Forest model driven by "distant context probes" achieved extraordinary performance levels. It successfully attained 97\% segmentation accuracy on distinguishing genuine play intent from background data. Furthermore, the model provided extremely tight temporal bounding, showing merely a 7\,ms average detection error on labeling constraints.

\section{Ablation Studies}

\subsection{Variance Entropy Scoring}
Considering that 80--90\% of raw game session data operates purely as noisy background idleness, attempting to run models without first applying the variance-based entropy score trimming produces severe data imbalance, crippling model convergence and execution speeds.

\subsection{Contextual Probes}
Ablating the "distant context probes" verifies that analyzing isolated instant frames limits classification to naive static thresholding. Including pre- and post-frame analytical slices serves to contextualize the momentum of a single finger, allowing 97\% classification precision.

\section{Parameter Optimizer Results}
The CMA-ES optimizer effectively simulated the hardware switch logic, validating the system's ability to hunt specific hardware tolerances to uniquely match physical telemetry. It successfully discovers configurations that improve objective measures by collapsing latency gaps within the parameters of strong recall and precision.

\section{End-to-End User Validation Protocol (Project ARROW)}
To close the loop from theoretical model accuracy to practical HCI outcomes, a structured end-to-end user validation protocol (Project ARROW) is designed. This evaluation assesses the subjective trust users place in the calibration recommendations and quantifies objective performance changes during actual gameplay. 

\subsection{Study Objectives}
The user validation protocol operates on three primary objectives:
\begin{enumerate}
    \item \textbf{Subjective Trust \& Perception:} Assessing users' willingness to trust an autonomous calibration tool and accept its recommendation, specifically tracking feedback regarding how the "intent proficiency score" and recommendation UI are interpreted.
    \item \textbf{Objective Performance (Zapper Trajectory):} Utilizing a custom standalone aiming minigame (Zapper) to analyze "stop-to-click" latency. A measurable reduction in this mechanical threshold latency verifies the optimizer's effectiveness over human-adjusted defaults.
    \item \textbf{Misclick Monitoring (Performance Loss):} Actively logging any catastrophic regressions, explicitly tracking if a recommended configuration leads to unintended inputs or if the user actively rejects the calibration due to lack of trust.
\end{enumerate}
(Placebo testing was intentionally excluded from the scope due to participant sample size limitations).

\subsection{Protocol Structure}
The 60-minute evaluation protocol is divided into four distinct phases:
\begin{itemize}
    \item \textbf{Phase 1: Pre-Session Onboarding (0--5 mins).} Users define their base ergonomic settings (DPI). The mouse is strictly reset to a conservative standard configuration baseline. Initial "out-of-the-box" impressions are recorded and a participant ID is assigned.
    \item \textbf{Phase 2: Baseline Calibration (5--20 mins).} Utilizing a think-aloud protocol, users engage with the initial calibration wizard and play for 10--15 minutes while telemetry is recorded. This phase concludes with a 1-minute baseline performance round in the Zapper minigame.
    \item \textbf{Phase 3: Post-Optimization UX (20--40 mins).} Users review the UI presenting the optimization recommendation. If a user rejects the algorithm's configuration due to a lack of trust, it is recorded as a critical UX failure alongside their justification. Users then apply the optimization and play a secondary session, culminating in a final 1-minute Zapper performance test to capture delta-latency improvements.
    \item \textbf{Phase 4: Offboarding \& Validation (50--60 mins).} Technical verification ensures the GSense tracking history is correctly encapsulated for the specific participant ID, followed by final qualitative user surveying and debriefing. 
\end{itemize}

\section{Discussion}
\subsection{Strengths}
The solution successfully operationalizes continuous hardware capability with machine learning inference, addressing the known difficulty of correctly tuning analog actuation devices manually. Identifying 1\,kHz intent with 7\,ms precision empowers the system to drastically override default performance blocks efficiently.

\subsection{User Trust and Adoption Roadblocks}
User evaluation provided a stark insight into adoption: users categorically distrust entirely autonomous hardware modification engines acting as black boxes. For production scale viability (such as integration into GHub and GHabits platforms), visual performance validations alongside objective data transparency (graphs outlining the recommended change) are structurally essential to foster adoption.
